\documentclass[12pt]{article}
\usepackage{amsmath,amssymb}
\usepackage{geometry}
\geometry{margin=1in}
\usepackage[hidelinks]{hyperref}
\setlength{\parindent}{0pt}
\usepackage{xcolor}
\usepackage{etoolbox}  %Supports toggle commands
\definecolor{darkblue}{RGB}{0,0,139}

%SETTINGS
\newtoggle{solutions} \togglefalse{solutions}

\begin{document}

\begin{center}
\textbf{Week 3 - Solution class exercise} \\
Introduction to Health Economics: EC 387 \\
Boston University \\
Spring 2026 \\
Professor Pauline Mourot
\end{center}

\vspace{1em}

The consumer has \$20,000 in income and her utility function is $u(x) = \sqrt(x)$.\\
She faces health risks such that
\begin{itemize}
    \item[-] She has a 50\% chance of having no medical expenses
    \item[-] She has a 40\% chance of having \$1,000 in medical expenses (mild illness)
    \item[-] She has a 10\% chance of having \$10,000 in medical expenses (severe illness)
\end{itemize}
\vspace{3mm}

She is considering buying a health insurance plan with
\begin{itemize}
    \item[-] A premium of \$1,000
    \item[-] A deductible of \$500
    \item[-] A coinsurance rate of 20\%
\end{itemize}

\textbf{Questions}:
\begin{enumerate}
    \item What is the consumer's expected out-of-pocket cost?

    \vspace{3mm}
    { \color{darkblue}

    The expected out-of-pocket cost is
    \begin{align*}
    \mathbb{E}(\text{OOP})
    &= 0.5 \times \$20{,}000 + 0.4 \times \$1{,}000 + 0.1 \times \$10,000\\
    &= \$1{,}400.
    \end{align*}
    The expected out-of-pocket cost when the individual is uninsured is \$1,400.
    }

    \item What is the expected cost to the insurer?

    \vspace{3mm}
    { \color{darkblue}

    The expected cost to the insurer is
    \begin{align*}
    \mathbb{E}(\text{Cost}_{insurer})
    &= 0.5 \times \$0 + 0.4 \times \left( (\$1{,}000 - \$500) \times 0.8 \right) + 0.1 \times \left( (\$10{,}000 - \$500) \times 0.8 \right) \\
    &= \$160 + \$760 \\
    &= \$920.
    \end{align*}
    The expected cost to the insurer is \$920.
    }

    \item What is the expected profit to the insurer?

    \vspace{3mm}
    { \color{darkblue}
    The expected profit to the insurer is
    \begin{align*}
    \mathbb{E}(\text{Profit}_{insurer})
    &= \text{Premium} - \mathbb{E}(\text{Cost}_{insurer}) \\
    &= \$1{,}000 - \$920 \\
    &= \$80.
    \end{align*}
    The expected profit to the insurer is \$80.
    }

    \item What would be the actuarially fair price for this insurance plan?
    \textit{Hint: the actuarially fair price is the premium that would give the insurer zero expected profit.}\\

    \vspace{3mm}
    { \color{darkblue}
    The actuarially fair price for this insurance plan is equal to the expected cost to the insurer,
    which we calculated in question 2:
    \begin{align*}
    \text{Actuarially fair price} = \mathbb{E}(\text{Cost}_{insurer}) = \$920.
    \end{align*}
    }

    \item Does the consumer prefer to remain uninsured or to buy this insurance plan?

    \vspace{3mm}
    { \color{darkblue}
    To determine the consumer's preference, we need to calculate the expected utility in both scenarios.
    \begin{itemize}
        \item Uninsured:
        \begin{align*}
        \mathbb{E}(U_{uninsured}) &= 0.5 \times \sqrt{20{,}000} + 0.4 \times \sqrt{20{,}000 - 1,000} + 0.1 \times \sqrt{20{,}000 - 10,000} \\
        &\approx 135.8469.
        \end{align*}

        \item Insured:
        \begin{align*}
        \mathbb{E}(U_{insured}) &= 0.5 \times \sqrt{20{,}000 - \underbrace{1,000}_{\text{premium}}} \\
        & + 0.4 \times \sqrt{20,000 - 500 - (1,000-500) \times 0.2 - \underbrace{1,000}_{\text{premium}}} \\
        &+ 0.1 \times \sqrt{20{,}000 - 500 - (10,000-500) \times 0,2 - \underbrace{1,000}_{\text{premium}}} \\
        &\approx 136.06.
        \end{align*}
    \end{itemize}
    Since $\mathbb{E}(U_{insured}) > \mathbb{E}(U_{uninsured})$, the consumer prefers to buy this insurance plan.
    }

    \item What is the maximum premium the consumer would be willing to pay for this insurance plan?

    \vspace{3mm}
    { \color{darkblue}
    Let $P_{max}$ be the maximum premium the consumer is willing to pay.
    We need to solve for $P_{max}$ such that the expected utility when insured
    equals the expected utility when uninsured:
    \begin{align*}
    & \mathbb{E}(U_{insured, P_{max}}) = \mathbb{E}(U_{uninsured})\\
    \leftrightarrow &0.5 \times \sqrt{20{,} - P_{max}} + 0.4 \times \sqrt{20,000 - 500 - (1000-500) \times 0.2 - P_{max}} \\
    &\text{  }+ 0.1 \times \sqrt{20{,}000 - 500 - (10,000-500) \times 0,2 - P_{max}} \approx 135.8469\\
    \end{align*}

    You can solve this equation numerically using Wolfram alpha (\url{https://www.wolframalpha.com/input}) or statistical softwares like R.
    Solving this equation numerically, we find that $P_{max} \approx 1,058.74$.
    }

\end{enumerate}




\end{document}